
\documentclass[11pt,a4paper]{article}
\usepackage[margin=20mm]{geometry}
\usepackage{kotex}
\usepackage{amsmath,amssymb,array,longtable,booktabs, calc}
\usepackage{xurl}
\usepackage{xcolor}
\usepackage{comment}
\usepackage[unicode, hidelinks]{hyperref}
\setlength{\parindent}{0pt}
\setlength{\parskip}{6pt}
\setlength{\emergencystretch}{3em}
\renewcommand{\arraystretch}{1.18}
\allowdisplaybreaks
\newcolumntype{P}[1]{>{\raggedright\arraybackslash}p{#1}}
\newcommand{\blocktitle}[1]{\par\bigskip\noindent\textbf{\large #1}\par\smallskip}

\setcounter{secnumdepth}{0}
\providecommand{\tightlist}{\setlength{\itemsep}{0pt}\setlength{\parskip}{0pt}}


\begin{document}
\blocktitle{④ 전체 시장과 multiple penalties 고려한 모델}

\[
\textbf{전체 시장 \(\rightarrow\) 로봇택시 선택수요 \(\rightarrow\) 차량 배치 \(\rightarrow\) 부족 이후 전환}
\]

\blocktitle{1. 수요를 두 종류로 만든다}

\[
D_c^{R,s}=M_c^s p_c^{R,s},
\qquad
D_c^{N,s}=M_c^s(1-p_c^{R,s})
\]

\begin{longtable}{@{}P{0.14\textwidth}P{0.41\textwidth}P{0.37\textwidth}@{}}
\toprule
기호 & 의미 & 주의점 \\
\midrule
\endhead
$M_c^s$ & 1년 동안 도시 $c$에서 발생한 모든 이동 수단에 걸친 잠재적 이동 기회의 총합 & 로봇수요 + 인간수요 + 기타수요 (ex. Uber/Lyft 완료 건수와 동일하지 않음) \\
$p_c^{R,s}$ & 배치 차량 수(용량) 제약을 적용하기 전, 지정된 서비스 조건 하에서 도시 $c$의 이동 기회에 대해 우리 회사 로봇택시 서비스를 선택할 확률 & \textcolor{red}{요금·대기시간 등 기준 서비스 조건을 명시} \\
$D_c^{R,s}$ & 로봇택시 선택수요 & 차량 부족을 적용하기 \textbf{이전} 수요 \\
$D_c^{N,s}$ & 로봇택시 비선택 수요 & 인간택시·다른 수단·이용 포기 등 \\
\bottomrule
\end{longtable}

\textbf{분모를 명확히 해야 해.} 기존의 \emph{완료된} 전체 ride-hailing 건수를 $M$으로 사용하면, 비선택 집단을 도보·이용 포기까지 포함한다고 해석할 수 없어. 이를 포함하려면 \textbf{잠재적인 이용 기회}를 추정해야 하고, 기존 ride-hailing 건수는 그 추정에 사용하는 대리자료가 돼.

\textcolor{red}{위 곱셈식은 \textbf{시나리오별 평균 선택량을 사용하는 근사}야. 개인 선택의 추가 변동까지 넣고, 조건부 독립·동일 선택확률을 가정한다면 다음으로 대체할 수 있어.}

※ M, p는 불확실한 변수.

{\color{red}\[\begin{aligned}
D_c^{R,s}&\sim\operatorname{Binomial}(M_c^s,p_c^{R,s}),\\
D_c^{N,s}&=M_c^s-D_c^{R,s}.
\end{aligned}
\]}

%%%%%%%%%%%%%%%%%%%%%%%%%%%%%%%%%%%%%%%%%%%%%%%%%%%%%%%%%%%%%%%%
\blocktitle{2. 핵심 모형: 도시 선택 + 차량 배치 + multiple penalties}

\[
\boxed{
\begin{aligned}
\min_{y,q,R,O,U,H,L}\quad
&\underbrace{\sum_c F_cy_c}_{\text{진입비용}}
+\underbrace{\sum_c K_cq_c}_{\text{차량비용}}\\
&+\frac1{|\mathcal S|}\sum_s\sum_c
\left[
\underbrace{c_c^+O_c^s}_{\text{남는 용량}}
+\underbrace{c_{c,H}^-H_c^s}_{\text{인간택시 전환 손실}}
+\underbrace{c_{c,L}^-L_c^s}_{\text{나머지 미처리 손실}}
\right].
\end{aligned}
}
\]

모든 수요·처리량·비용은 \textbf{같은 평가기간} 기준이야. \textcolor{red}{※ 평가기간은 연구자가 정하고, 모델은 그 기간에 적용할 도시/차량 배치 결정한다..} \textbf{$c^+,c^-$에서 $+,-$는 비용의 부호가 아니라 과잉·부족의 구분}이야.


\blocktitle{3. 모형에 사용하는 기호}
%공통기호
\begin{longtable}{@{}P{0.12\textwidth}P{0.34\textwidth}P{0.46\textwidth}@{}}
\toprule
기호 & 의미 & 설정방법 \\
\midrule
\endhead
$\mathcal C$ & 잠재 후보 도시 집합 & 회사의 후보 목록. 현재 회사 현황 자료는 참고 \\
$\mathcal S$ & 수요 시나리오 집합 & 연구자가 생성 \\
$y_c$ & 도시 $c$에 진입하면 1 & \textbf{최적화가 결정} \\
$q_c$ & 도시 $c$의 배치 차량 수 & \textbf{최적화가 결정} \\
$\kappa_c$ & 차량 한 대의 기간 내 처리 가능 trip 수 & \textbf{운영자료 또는 가정}. 처음에는 고정 \\
$Q$ & 전체 배치 가능한 차량 수 & 회사 입력. 없으면 연구자가 실험값 설정 \\
$\bar q_c$ & 도시별 차량 배치 상한 & 시설·운영 한도. 별도 한도가 없으면 $Q$ \\
$a_c$ & 계획 시점에 해당 도시에서 운영 가능하면 1 & 허가·ODD·시설자료로 판정. 미확인은 별도 가정 \\
$B$ & 동시에 진입할 수 있는 도시 수의 상한 & \textbf{실제 제한이 있을 때만} 회사 입력 또는 실험값 \\
\bottomrule
\end{longtable}

%수요변수
\begin{longtable}{@{}P{0.16\textwidth}P{0.42\textwidth}P{0.34\textwidth}@{}}
\toprule
결과변수 & 의미 & 설정 방법 \\
\midrule
\endhead
$M_c^s$ & 1년 동안 도시 $c$에서 발생한 모든 이동 수단에 걸친 잠재적 이동 기회의 총합 & \textcolor{red}{(입력값) 추정}\\
$p_c^{R,s}$ & 배치 차량 수(용량) 제약을 적용하기 전, 지정된 서비스 조건 하에서 도시 $c$의 이동 기회에 대해 우리 회사 로봇택시 서비스를 선택할 확률 & \textcolor{red}{(입력값) 추정}\\
$D_c^{R,s}$ & 로봇택시 선택수요 & 주어진 $M_c^s, p_c^{R,s}$로 계산 \\
$D_c^{N,s}$ & 로봇택시 비선택 수요(인간택시·다른 수단·이용 포기 등) & 주어진 $M_c^s, p_c^{R,s}$로 계산 \\
$\tau_c^s$ & 미처리 고객 중 인간택시로 전환 비율 & \\
\bottomrule
\end{longtable}
\textbf{※ 수요 관련 변수는 입력값인 $p_c^{R,s}, M_c^s$와 모델에서 설정한 배치 및 전환조건에 따라 확률적으로 결정되는 결과변수.}

%결과변수
\begin{longtable}{@{}P{0.16\textwidth}P{0.42\textwidth}P{0.34\textwidth}@{}}
\toprule
결과변수 & 의미 & 설정 방법 \\
\midrule
\endhead
$R_c^s$ & 실제 로봇택시 처리량 & 모형이 배치결과에 따라 계산 \\
$O_c^s$ & 사용하지 못한 처리용량 & 모형이 배치결과에 따라 계산 \\
$U_c^s$ & 로봇택시 미처리 수요 & 모형이 배치결과에 따라 계산 \\
$H_c^s$ & 미처리 후 인간택시로 전환되어 처리된 수요 & 아래전환 조건으로 계산($H_c^s=\tau_c^sU_c^s$) \\
$L_c^s$ & 나머지 미처리 수요 & 다른 수단·이용 포기 등을 묶음 ($L_c^s=(1-\tau_c^s)U_c^s$) \\
\bottomrule
\end{longtable}

\textbf{※ 결과변수는 각각 데이터를 조사해서 예측하는 대상이 아니다. 수요·배치·전환 조건으로 계산되는 결과.}

%%%%%%%%%%%%%%%%%%%%%%%%%%%%%%%%%%%%%%%%%%%%%%%
\blocktitle{4. 공통 제약식}

\textbf{A. Definition} \\
\textbf{1)수요 $D_c^s$와 공급능력 $\kappa_cq_c$가 결정되면 배치 결과에 따라 계산되는 값}은 다음 세 가지로, 별도로 조사할 변수가 아니다. 여기서 $(x)^+=\max\{x,0\}$. 세 값의 단위는 모두 \textbf{trip 수}.

\[
\begin{aligned}
R_c^s&=\min\{D_c^s,\kappa_cq_c\}
&&\text{실제 처리량},\\
O_c^s&=(\kappa_cq_c-D_c^s)^+
&&\text{남는 처리용량},\\
U_c^s&=(D_c^s-\kappa_cq_c)^+
&&\text{미처리 잠재수요}.
\end{aligned}
\]

\textbf{2)Human fallback은 전환비율 $\tau_c^s$을 입력으로 설정한다.}
\[
H_c^s=\tau_c^sU_c^s,\qquad
L_c^s=(1-\tau_c^s)U_c^s
\]
$\tau_c^s\in[0,1]$는 \textbf{미처리 수요 중 실제로 인간택시로 전환되어 처리되는 비율}이야.

이 방식은 간단하지만, 비용항을 정리하면

\[
c_{c,H}^-H_c^s+c_{c,L}^-L_c^s
=
\underbrace{
\left[\tau_c^sc_{c,H}^-+(1-\tau_c^s)c_{c,L}^-\right]
}_{\text{유효 부족비용}}
U_c^s.
\]

\textbf{즉, 전환비율이 고정이면 multiple penalties는 하나의 가중 부족비용으로 축약된다.}

\[\boxed{
\begin{aligned}
y_c\le q_c&\le\bar q_cy_c
&&\forall c,\\
\sum_cq_c&\le Q,\\
y_c&\le a_c
&&\forall c,\\[2pt]
R_c^s+O_c^s&=\kappa_c^sq_c
&&\forall c,s,\\
R_c^s+U_c^s&=D_c^{R,s}
&&\forall c,s,\\
U_c^s&=H_c^s+L_c^s
&&\forall c,s,\\
y_c&\in\{0,1\},\quad q_c\in\mathbb Z_{\ge0},\\
R_c^s,O_c^s,U_c^s,H_c^s,L_c^s&\ge0.
\end{aligned}
}\]

\begin{longtable}{@{}P{0.32\textwidth}P{0.62\textwidth}@{}}
\toprule
제약 & 의미 \\
\midrule
\endhead
$y_c\le q_c\le\bar q_cy_c$ & 진입하지 않으면 0대, 진입하면 1대 이상 \\
$\sum_cq_c\le Q$ & 총 차량 제한 \\
$y_c\le a_c$ & 계획 시점에 운영 가능한 도시에만 진입 \\
$R_c^s+O_c^s=\kappa_c^s q_c$ & 전체 처리용량 = 사용한 용량 + 남는 용량 \\
$R_c^s+U_c^s=D_c^{R,s}$ & 로봇택시 수요 = 처리량 + 미처리량 \\
$U_c^s=H_c^s+L_c^s$ & 미처리 수요의 결과를 두 종류로 분해 \\
\bottomrule
\end{longtable}

%%%%%%%%%%%%%%%%%%%%%%%%%%%%%%%%%%%%%%%%%%%%%%%%%%%%%%%%%%%%%%%%%
\blocktitle{※ 입력값은 무엇으로 설정하는가?}

\textbf{`자료로 추정', \textcolor{red}{`회사 입력'}, `연구자 가정'을 구분해서 표시하면 돼.} 시나리오 입력도 생성 후 최적화에서는 고정값이야.

\begin{longtable}{@{}P{0.14\textwidth}P{0.24\textwidth}P{0.24\textwidth}P{0.28\textwidth}@{}}
\toprule
입력 & 설정 방법 & 기존 조사자료 활용 & 아직 필요한 자료 \\
\midrule
\endhead
$M_c^s$ & \textbf{자료 추정 + 시장 정의 가정} & 인구·고용·기존 ride-hailing·공항·상업활동 & 잠재 이동시장 분모. 완료 trip만으로는 부족 \\
$p_c^{R,s}$ & \textbf{선택자료 추정}, 없으면 가정 & 소득·무차량 가구·교통 접근성 등 후보 설명변수 & 로봇택시 선택·비선택, 제시 요금·대기시간·대안 \\
$\kappa_c^s$ & \textbf{운영자료 추정}, 없으면 가정 & 혼잡·도로·기상·충전 여건 & 운행·픽업·재배치·충전·정비시간 \\
$\tau_c^s$ & \textbf{전환자료 추정}, 없으면 가정 & 현재 자료로 직접 추정 불가 & 로봇택시 미처리 이후 실제 행동 \\
$a_c$ & \textbf{자료로 판정 + 개시 조건 가정} & 허가·ODD·거점 조건 & 해당 회사·계획 시점의 조건 \\
$\bar q_c$ & \textbf{시설·운영자료 또는 회사 입력} & 충전·주차·정비 여건 & 실제 확보 가능 규모. 별도 상한 없으면 $Q$ \\
$Q$ & \textbf{회사 입력 또는 실험값} & 수요자료로 추정하지 않음 & 가용 차량 수 \\
$F_c$ & \textbf{추가 비용자료 또는 가정} & 거점·부지 비용 일부 & 초기 구축·진입 준비비 \\
$K_c$ & \textbf{추가 비용자료 또는 가정} & 현재 조사로 확정 불가 & 평가기간 차량 배치·보유비 \\
$c_c^+$ & \textbf{회사 손실자료 또는 가정} & 현재 조사로 확정 불가 & 잔여용량 부담. $K_c$와 중복 금지 \\
$c_{c,H}^-,c_{c,L}^-$ & \textbf{회사 손실자료 또는 가정} & 현재 조사로 확정 불가 & 전환 유형별 기회손실·고객 이탈 차이 \\
\bottomrule
\end{longtable}

$y_c,q_c$는 이 표에 넣지 않아. \textbf{연구자가 설정하는 입력이 아니라 모형이 결정하는 값}이기 때문이야.

%%%%%%%%%%%%%%%%%%%%%driving factors%%%%%%%%%%%%%%%%%%%%%%%%%%%%
\blocktitle{5. Driving factors}

\hypertarget{muxc758-driving-factors-uxadfcuxac70-uxb17cuxbb38-uxcd94uxac00}{%
\subsection{\texorpdfstring{1. \(M\)의 driving factors}{A. M의 driving factors}}\label{muxc758-driving-factors-uxadfcuxac70-uxb17cuxbb38-uxcd94uxac00}}

아래 논문들은 주로 \textbf{관측된 ride-hailing 이용량}을 설명하기 때문에 M의 정의에 포함되는 '미실현 수요'까지 포함되지 않아 직접 추정은 불가능. 하지만 논문들이 제공하는 수요 패턴은 '잠재 이동 기회 \(M\)'을 모형적으로 구성하는데 중요한 단서가 된다.

\begin{longtable}[]{@{}
  >{\raggedright\arraybackslash}p{(\columnwidth - 6\tabcolsep) * \real{0.1600}}
  >{\raggedright\arraybackslash}p{(\columnwidth - 6\tabcolsep) * \real{0.3200}}
  >{\raggedright\arraybackslash}p{(\columnwidth - 6\tabcolsep) * \real{0.3200}}
  >{\raggedright\arraybackslash}p{(\columnwidth - 6\tabcolsep) * \real{0.2000}}@{}}

  
\toprule\noalign{}
\begin{minipage}[b]{\linewidth}\raggedright
변수
\end{minipage} & \begin{minipage}[b]{\linewidth}\raggedright
근거 논문
\end{minipage} & \begin{minipage}[b]{\linewidth}\raggedright
논문에서 확인한 범위
\end{minipage} & \begin{minipage}[b]{\linewidth}\raggedright
\textcolor{red}{측정자료}
\end{minipage} \\


\midrule\noalign{}
\endhead


\bottomrule\noalign{}
\endlastfoot

\textbf{인구·인구밀도} & Elkhouly·Alhadidi, \emph{Ride-sharing
Determinants} (\href{https://arxiv.org/abs/2406.11590}{논문}); Zhang 외,
\emph{Examining spatial heterogeneity\ldots{}}
(\href{https://arxiv.org/abs/2209.07980}{논문}) & 인구 및 지역 밀도와
일반 ride-hailing 이용량의 관계. 밀도 효과는 지역에 따라 달라질 수 있음
& \href{https://www.census.gov/programs-surveys/acs}{ACS} \\

\textbf{고용·경제활동} & Elkhouly·Alhadidi
(\href{https://arxiv.org/abs/2406.11590}{논문}) & 경제활동인구·고용환경을 이용량
설명에 사용 & \href{https://lehd.ces.census.gov/data/}{LEHD/LODES},
ACS \\

\textbf{소득} & Elkhouly·Alhadidi
(\href{https://arxiv.org/abs/2406.11590}{논문}) & Chicago에서 소득과
이용량의 관계 분석 &
\href{https://data.census.gov/table/ACSDT5Y2024.B19013}{ACS B19013} \\

\textbf{무차량 가구 비율} & Elkhouly·Alhadidi
(\href{https://arxiv.org/abs/2406.11590}{논문}) & 자동차를 보유하지 않은
가구와 이용량의 관계 분석 &
\href{https://data.census.gov/table/ACSDT1Y2022.B08201}{ACS B08201} \\

\textbf{토지이용 혼합·접근성·상업활동} & Zhang 외, \emph{Examining
spatial heterogeneity\ldots{}}
(\href{https://arxiv.org/abs/2209.07980}{논문}) & 도시환경과
ridesourcing 이용량의 비선형·공간적 차이 &
\href{https://www.epa.gov/smartgrowth/smart-location-mapping}{EPA
SLD} \\

\textbf{식당·여가시설 밀도} & Loa·Lee·Circella, \emph{Just another
ride-hailing?}
(\href{https://link.springer.com/article/10.1007/s11116-026-10751-9}{논문})
& \textbf{SF 로봇택시 이용 건수}와의 관계. \(M\)과 \(p\)의 효과를 분리한
결과는 아님 &
\href{https://www.census.gov/programs-surveys/cbp.html}{Census CBP},
OpenStreetMap \\

\textbf{대중교통·보행환경} & Zhang 외 
(\href{https://arxiv.org/abs/2209.07980}{논문}); Cats 외, \emph{Beyond
the dichotomy\ldots{}}
(\href{https://doi.org/10.1371/journal.pone.0262496}{논문}) & 지역환경
및 대중교통과의 시간 경쟁력·대체 가능성 분석. 일률적인 증가·감소 효과는
아님 & EPA SLD, \href{https://gtfs.org/documentation/overview/}{GTFS} \\

\textbf{공항 관련 시장규모} & Zhang 외
(\href{https://arxiv.org/abs/2209.07980}{논문}) & 공항 이동을 별도
공간적 유형으로 분석. \textbf{공항 여객 수의 직접 효과를 추정한 것은
아님} &
\href{https://www.faa.gov/airports/planning_capacity/passenger_allcargo_stats/passenger}{FAA
여객자료}. 시장규모 대리변수로 사용 \\

\textbf{시간·요일·계절} & Elkhouly·Alhadidi
(\href{https://arxiv.org/abs/2406.11590}{논문}) & 시간에 따른 이용량
차이 분석. 휴일 등 구체적 변수 선택은 별도 설계 & Trip 자료의
날짜·시간 \\

\textbf{기상} & \textbf{추가:} Liu·Jiang·Chen, \emph{Quantifying the
impact of weather on ride-hailing ridership: Evidence from Haikou,
China} (2021, \href{https://doi.org/10.1016/j.tbs.2021.04.002}{논문}) &
날씨와 ride-hailing 이용량의 관계 분석. 미국 도시에서의 효과는 재검증
필요 & \href{https://www.ncei.noaa.gov/cdo-web/}{NOAA CDO} \\
\end{longtable}

$M$의 결정요인은 도시특성(원안) + 시간/계절 맥락(상황) + 시장 활동(결과: 과거 Uber/Lyft 이용량)\textbf{과거 이용량}은 별도로 사용할 수 있는 예측변수야. 다만 이것은 ''왜 수요가 발생하는가''를 설명하지 못하지, \textbf{수요의 크기(규모)를 보여주는 결과 데이터}로 구분하면 돼.


%%%%%%%%%%%%%%%%p의 driving factors%%%%%%%%%%%%%%%%%%%%%%%%%
%%%%%%%%%%%%%%%%%%%%%%%%%%%%%%%%%%%%%%%%%%%%%%%%%%%%%%%%%%%%
\hypertarget{puxc758-driving-factors-uxadfcuxac70-uxb17cuxbb38-uxcd94uxac00}{%
\subsection{\texorpdfstring{2. \(p\)의 driving factors}{B. p의 driving factors}}\label{puxc758-driving-factors-uxadfcuxac70-uxb17cuxbb38-uxcd94uxac00}}

\begin{longtable}[]{@{}
  >{\raggedright\arraybackslash}p{(\columnwidth - 6\tabcolsep) * \real{0.1600}}
  >{\raggedright\arraybackslash}p{(\columnwidth - 6\tabcolsep) * \real{0.3200}}
  >{\raggedright\arraybackslash}p{(\columnwidth - 6\tabcolsep) * \real{0.3000}}
  >{\raggedright\arraybackslash}p{(\columnwidth - 6\tabcolsep) * \real{0.2200}}@{}}

  
\toprule\noalign{}
\begin{minipage}[b]{\linewidth}\raggedright
변수
\end{minipage} & \begin{minipage}[b]{\linewidth}\raggedright
근거 논문
\end{minipage} & \begin{minipage}[b]{\linewidth}\raggedright
적용 범위
\end{minipage} & \begin{minipage}[b]{\linewidth}\raggedright
필요한 측정자료
\end{minipage} \\


\midrule\noalign{}
\endhead


\bottomrule\noalign{}
\endlastfoot

\textbf{상대 요금} & \textbf{추가:} Yin·Cherchi, \emph{Willingness to
pay for automated taxis\ldots{}} (2024; 온라인 2022,
\href{https://doi.org/10.1007/s11116-022-10319-3}{논문}) & 일반택시와
자율주행 택시의 선택실험 & 양쪽 제시요금 + 선택결과 \\

\textbf{상대 대기시간} & Yin·Cherchi
(\href{https://doi.org/10.1007/s11116-022-10319-3}{논문}) & 서비스 조건에 따른 선택을 분석 & 양쪽 제시 대기시간 + 선택결과 \\

\textbf{차내 이동시간} & Yin·Cherchi
(\href{https://doi.org/10.1007/s11116-022-10319-3}{논문}) & 요금·시간 사이의 선택상충 관계 분석
& 양쪽 제시 이동시간 \\

\textbf{신뢰·가격가치·이용습관} & \textbf{추가:} Wei 외, \emph{Using the
Extended Unified Theory\ldots{}} (2024,
\href{https://www.nature.com/articles/s41599-024-03271-3}{논문}) &
로봇택시 수용의향·행동 연구. \textbf{설문 척도를 바로 선택확률로 해석할
수 없음} & 개인 태도·습관 설문 \\

\textbf{연령·로봇택시 관련 지식} & Yin·Cherchi
(\href{https://doi.org/10.1007/s11116-022-10319-3}{논문}) & 개인 특성에
따른 선택 차이. ``젊을수록 무조건 선호'' 같은 부호 고정은 부적절 & 개인
특성과 선택자료 연결 \\

\textbf{이용후기·차내 지원기능} & Yin·Cherchi
(\href{https://doi.org/10.1007/s11116-022-10319-3}{논문}) & 고객평가,
목적지 변경·운영자 소통 등의 기능을 선택속성으로 분석 & 선택실험의
서비스 속성 \\

\textbf{로봇택시 이용 경험} & \textbf{추가:} Yin·Cherchi·Nettle,
\emph{Living in a City Where Automated Taxis are Operating and Using
Them\ldots{}} (2023,
\href{https://doi.org/10.1177/03611981231155181}{논문}) & 거주지역의
운행 여부와 개인 이용 경험에 따른 선호 차이 분석 & 개인 이용 경험 \\

\textbf{대중교통 등 대안의 경쟁력} & Cats 외, \emph{Beyond the dichotomy\ldots{}}
(\href{https://doi.org/10.1371/journal.pone.0262496}{논문}) &
\textbf{일반 ride-hailing의 근거}. 로봇택시 고유의 선호효과는 아님 &
대안의 이용 가능 여부·요금·시간 \\
\end{longtable}

앞서 제시한 \textbf{픽업 도보시간·동행자·소득·차량 보유}는 여전히 검토할
후보야. 하지만 이번에 확인한 자료만으로 각각을 ``로봇택시 선택에 유의한
효과가 입증된 변수''라고 표시하지는 않겠어.



%%%%%%%%%%%%%%%%%%%%%%%%%%%%%%%%%%%%%%%%%%%%%%%%%%%%%%%%%%%%%%%%%%%%%%%%%%%%%%%%%%%%%%%%%%%%%%%%%%%%%%%%%%%%%%%%%%%%%%%%%%%%%%%%%%%%%%%%%%%
\begin{comment}

\blocktitle{교수님이 multiple penalties를 요구하면: ②의 부족손실을 확장}

\[
c_c^-U_c^s
\quad\longrightarrow\quad
c_{c,H}^-H_c^s+c_{c,L}^-L_c^s
\]

\[
H_c^s=\tau_c^sU_c^s,\qquad
L_c^s=(1-\tau_c^s)U_c^s.
\]

\begin{longtable}{@{}P{0.20\textwidth}P{0.32\textwidth}P{0.40\textwidth}@{}}
\toprule
기호 & 의미 & 설정 방법 \\
\midrule
\endhead
$\tau_c^s$ & 미처리 고객 중 인간택시로 전환하는 비율 & \textbf{추가 선택·전환조사 필요}, 없으면 가정 \\
$H_c^s$ & 인간택시 전환 수요 & 위 비율로 계산 \\
$L_c^s$ & 나머지 전환·이용 포기 수요 & 위 비율로 계산 \\
$c_{c,H}^-,c_{c,L}^-$ & 두 경우의 회사 손실 & 추가자료 또는 가정. 어느 쪽이 더 큰지는 근거 필요 \\
\bottomrule
\end{longtable}

현재 확보한 CPUC 취소 사유에는 \textbf{취소 후 어떤 수단을 택했는지}가 없어. 따라서 이 확장은 우선 \textbf{가정 기반 실험}으로 제시해야 해.

\blocktitle{3. 네가 조사한 자료는 실제로 어디에 들어가나?}

\begin{longtable}{@{}P{0.23\textwidth}P{0.31\textwidth}P{0.38\textwidth}@{}}
\toprule
모형 입력 & 기존 자료의 활용 & 아직 필요한 것 \\
\midrule
\endhead
\textbf{$D_c^s$: 수요} & 인구·소득·고용·교통·POI 등을 이용량 예측에 연결 & 완료량과 잠재수요의 차이를 다룰 근거. 개인 선택확률을 추정하려면 선택·비선택 자료 \\
\textbf{$\kappa_c$: 차량 처리량} & 혼잡·도로·기상·충전 여건을 활용 & 탑승·픽업·재배치 시간, 운영·충전·정비시간 \\
\textbf{$a_c$: 운영 가능성} & 규제·허가·ODD·거점 조건 활용 & 회사별 기술 범위와 계획 시점의 조건 \\
\textbf{$\bar q_c$: 도시별 상한} & 거점·충전·주차 여건 활용 & 실제 확보 가능한 시설 용량 \\
\textbf{$F_c,K_c,m_c,c_c^\pm$: 경제적 입력} & 토지·거점 비용 자료 일부 참고 & 회사별 비용·수익 자료. 현재 조사로 대부분 확정 불가 \\
\textbf{$Q,B,\rho$: 자원·선호} & 수요자료로 추정하는 대상이 아님 & 회사가 주거나, 연구자가 실험 조건으로 명시 \\
\bottomrule
\end{longtable}

수요자료의 출발점은
\href{https://www.cpuc.ca.gov/regulatory-services/licensing/transportation-licensing-and-analysis-branch/autonomous-vehicle-programs/quarterly-reporting}{CPUC 실제 이용자료},
도시 설명변수는
\href{https://www.census.gov/data/developers/data-sets/acs-5year.html}{ACS}·\href{https://www.epa.gov/smartgrowth/smart-location-mapping}{EPA SLD}야.
\textbf{현재 확인한 공개 이용량만으로 잠재수요와 개인 선택확률까지 추정됐다고 할 수는 없어.}

\textbf{교수님께 가져갈 구성은 `①을 기본모형으로 제시 \(\rightarrow\) ②·③은 경제적 목적의 대안 \(\rightarrow\) multiple penalties는 자료 또는 가정에 따른 확장'이면 돼.} 세 모형 모두 수요 시나리오를 도시 정보에 조건부로 생성하면 contextual 모형이 되고, 목적함수를 별도로 복잡하게 바꿀 필요는 없어.




%%%%%%%%%%%%%%%%%%%%%%%%%%%%%%%%%%%%%%%%%%%%%%%%%%%%%%%%%%%%
%%%%%%%%%%%%%%%%%%%%%%%참고%%%%%%%%%%%%%%%%%%%%%%%%%%%%%%%%%%
============================================================
\blocktitle{3. Human fallback은 다음 두 방식 중 하나를 사용한다}

\textbf{A. 먼저 구현할 형태: 전환비율을 입력으로 설정}

\[
\boxed{
H_c^s=\tau_c^sU_c^s,\qquad
L_c^s=(1-\tau_c^s)U_c^s
}
\]

$\tau_c^s\in[0,1]$는 \textbf{미처리 수요 중 실제로 인간택시로 전환되어 처리되는 비율}이야.

이 방식은 간단하지만, 비용항을 정리하면

\[
c_{c,H}^-H_c^s+c_{c,L}^-L_c^s
=
\underbrace{
\left[\tau_c^sc_{c,H}^-+(1-\tau_c^s)c_{c,L}^-\right]
}_{\text{유효 부족비용}}
U_c^s.
\]

\textbf{즉, 전환비율이 고정이면 multiple penalties는 하나의 가중 부족비용으로 축약돼.} 이 점도 교수님께 설명하면 좋아.

\textbf{B. 공급 불확실성까지 반영: 인간택시의 추가 수용 여력 제한}

A의 등식을 \textbf{다음 조건으로 대체}해.

\[
\boxed{
\begin{aligned}
H_c^s&\le\tau_c^sU_c^s,\\
H_c^s&\le A_c^s,\\
L_c^s&=U_c^s-H_c^s
\end{aligned}
}
\]

\begin{longtable}{@{}P{0.20\textwidth}P{0.74\textwidth}@{}}
\toprule
추가 기호 & 의미 \\
\midrule
\endhead
$\tau_c^s$ & 미처리 고객 중 인간택시로 전환할 의향이 있는 비율 \\
$A_c^s$ & 인간택시가 추가로 처리할 수 있는 trip 수 \\
\bottomrule
\end{longtable}

이때 $c_{c,L}^->c_{c,H}^-\ge0$를 가정하면 모형은 가능한 만큼 인간택시로 전환해:

\[
H_c^s=\min\{\tau_c^sU_c^s,A_c^s\}.
\]

\textbf{이 비용 순서는 회사 손실에 대한 가정이지, 자동으로 성립하는 사실은 아니야.}

인간택시 기존 수요까지 명시하려면 추가로:

\[
\boxed{
\begin{aligned}
D_c^{H,s}&=\theta_c^sD_c^{N,s},\\
A_c^s&=(C_c^{H,s}-D_c^{H,s})^+
\end{aligned}
}
\]

\begin{longtable}{@{}P{0.20\textwidth}P{0.74\textwidth}@{}}
\toprule
기호 & 의미 \\
\midrule
\endhead
$\theta_c^s$ & 최초 로봇택시 비선택자 중 인간택시를 선택하는 비율 \\
$D_c^{H,s}$ & 처음부터 인간택시를 선택한 수요 \\
$C_c^{H,s}$ & 인간택시의 전체 처리용량 \\
\bottomrule
\end{longtable}

이는 \textbf{기존 인간택시 고객을 먼저 처리하고 남는 용량으로 fallback을 처리한다}는 근사야. 인간택시 공급의 불확실성은 $C_c^{H,s}$, 로봇택시 공급의 불확실성은 $\kappa_c^s$로 표현해.


%%%%%%%%%%%%%%%%%%%%%%%%%%%%%%%%%%%%%%%%%%%%%%%%%%%%%%%%%%%%%%%%%%%%%

\blocktitle{4. 입력값은 무엇으로 설정하는가?}

\textbf{`자료로 추정', `회사 입력', `연구자 가정'을 구분해서 표시하면 돼.} 시나리오 입력도 생성 후 최적화에서는 고정값이야.

\begin{longtable}{@{}P{0.14\textwidth}P{0.24\textwidth}P{0.24\textwidth}P{0.28\textwidth}@{}}
\toprule
입력 & 설정 방법 & 기존 조사자료 활용 & 아직 필요한 자료 \\
\midrule
\endhead
$M_c^s$ & \textbf{자료 추정 + 시장 정의 가정} & 인구·고용·기존 ride-hailing·공항·상업활동 & 잠재 이동시장 분모. 완료 trip만으로는 부족 \\
$p_c^{R,s}$ & \textbf{선택자료 추정}, 없으면 가정 & 소득·무차량 가구·교통 접근성 등 후보 설명변수 & 로봇택시 선택·비선택, 제시 요금·대기시간·대안 \\
$\kappa_c^s$ & \textbf{운영자료 추정}, 없으면 가정 & 혼잡·도로·기상·충전 여건 & 운행·픽업·재배치·충전·정비시간 \\
$\tau_c^s$ & \textbf{전환자료 추정}, 없으면 가정 & 현재 자료로 직접 추정 불가 & 로봇택시 미처리 이후 실제 행동 \\
$A_c^s$ & \textbf{운영자료 추정}, 없으면 가정 & 일반 ride-hailing 운영자료 일부 & 인간택시 가동시간·수요·잔여용량 \\
$\theta_c^s,C_c^{H,s}$ & B를 세분화할 때 추가 & TNC 자료·이동조사 일부 & 최초 인간택시 선택과 전체 공급능력 \\
$a_c$ & \textbf{자료로 판정 + 개시 조건 가정} & 허가·ODD·거점 조건 & 해당 회사·계획 시점의 조건 \\
$\bar q_c$ & \textbf{시설·운영자료 또는 회사 입력} & 충전·주차·정비 여건 & 실제 확보 가능 규모. 별도 상한 없으면 $Q$ \\
$Q$ & \textbf{회사 입력 또는 실험값} & 수요자료로 추정하지 않음 & 가용 차량 수 \\
$F_c$ & \textbf{추가 비용자료 또는 가정} & 거점·부지 비용 일부 & 초기 구축·진입 준비비 \\
$K_c$ & \textbf{추가 비용자료 또는 가정} & 현재 조사로 확정 불가 & 평가기간 차량 배치·보유비 \\
$c_c^+$ & \textbf{회사 손실자료 또는 가정} & 현재 조사로 확정 불가 & 잔여용량 부담. $K_c$와 중복 금지 \\
$c_{c,H}^-,c_{c,L}^-$ & \textbf{회사 손실자료 또는 가정} & 현재 조사로 확정 불가 & 전환 유형별 기회손실·고객 이탈 차이 \\
\bottomrule
\end{longtable}

$y_c,q_c$는 이 표에 넣지 않아. \textbf{연구자가 설정하는 입력이 아니라 모형이 결정하는 값}이기 때문이야.









%%%%%%%%%%%%%%%%%%%%%%%%%%%%%%%%%%%%%%%%%%%%%%%%%%%%%%%%%%%%%%%%%%%%

\blocktitle{5. Contextual은 어디에 들어가나?}

\[
\boxed{
X
\longrightarrow
\left\{
M_c^s,\ p_c^{R,s},\ \kappa_c^s,\ \tau_c^s,\ A_c^s
\right\}_{c,s}
\longrightarrow
\text{위 배치모형}
}
\]

이것들을 \textbf{처음부터 전부 예측해야 한다는 뜻은 아니야.}

\begin{longtable}{@{}P{0.23\textwidth}P{0.71\textwidth}@{}}
\toprule
단계 & 적용 \\
\midrule
\endhead
\textbf{첫 구현} & $X$ 없이 수요 시나리오를 주고, $\kappa,\tau$는 고정 \(\rightarrow\) 모형 A ($M$, $p$, $k$, $\tau$를 직접 세팅 \\
\textbf{Contextual 확장} & 도시·계절 정보로 수요 시나리오 생성 \\
\textbf{공급 확장} & 기상·운영정보로 $\kappa_c^s,A_c^s$ 변화 \(\rightarrow\) 모형 B \\
\textbf{가격·대기시간 확장} & 회사의 가격·차량 결정에 따라 선택확률이 달라지도록 확장 \\
\bottomrule
\end{longtable}

메모의 ``given $X$ \(\rightarrow\) scenario''는 조건부 예측과 잔차를 SAA로 연결하는 방향이고, 가격·차량 수가 수요 자체를 바꾸면 decision-dependent 확장이 돼.
\href{https://arxiv.org/abs/2207.13554}{Kannan 외},
\href{https://arxiv.org/abs/2406.20004}{Zhu·Yu·Bayraksan}

\textbf{교수님께는 위 목적함수·공통 제약식에 A를 붙인 모형을 기본안으로, B를 `인간택시 수요와 불확실한 공급까지 반영한 확장안'으로 제시하자.} 미진입 도시의 $U$는 실제 거절 고객이 아니라 \textbf{확보하지 못한 시장 기회}로 해석한다는 점도 함께 명시해야 해.


\end{comment}

%%%%%%%%%%%%%%%%%%%%%%%%%%%%%%%%%%%%%%%%%%%%%%%%%%%%%%%%%%%%%%%%%%%%%%%%%%%%%%%%%%%%%%%%%%%%%%%%%%%%%%%%%%%%%%%%%%%%%%%%%%%%%%%%%%%%%%%%%%%%%%%%%%%%%%%%%%%%%%%%%%%%%

\blocktitle{6. Things still need to be clarified}
1. how to set scenarios?($M$,$p$,$\tau$) \\
2. 기존자료 조사로 추정할 수 있는 것은? \\
3. 더 필요한 자료는? \\
4. 자료가 없는 기호에 대해서는 어떻게 가정을 해야하는가? \\
5. 목적함수 식이 해가 구해지는 것인가? \\
6. 추정해야 하는 결과변수 외 나머지 기호들에 대한 자료를 구하지 못할 경우 어떻게 세팅을 할 것인가? \\

\blocktitle{※ 무엇을 해야 하는가?}
1. 첫 구현의 범위를 고정: \(\kappa,\tau\)와 기준 요금·대기시간을 고정하고, 수요 시나리오를 넣어 현재 배치모형부터 실행한다. 가상 입력이면 ‘모형 검증용 실험’이라고 표시한다.\\
2. 자료와 수요를 연결: 같은 지역·기간의 \((X,M)\) 자료를 만들고, 별도로 로봇택시 선택확률 \(p\)을 추정할 자료가 있는지 확인한다. Kannan의 SAA 논문을 따라 잔차 기반 시나리오로 확장한다. \\
3. Driving factors를 결과로 확인: 경쟁 요금·대기시간·시장규모·처리능력을 바꿨을 때 선정 도시, 배치 대수, 미처리량, 총손실이 얼마나 변하는지 비교한다




\end{document}
